% ---------------------------------------------------------------------
%  CHƯƠNG 1 — MỞ ĐẦU
% ---------------------------------------------------------------------
\chapter{Mở đầu và cơ sở lý thuyết}
\label{ch:modau}

\begin{muctieu}
Chương này nêu bối cảnh, mục tiêu và phạm vi đồ án, sau đó trình bày cơ sở
lý thuyết cùng công nghệ của nền tảng AI ExamHub DAU.
\end{muctieu}

\section{Lý do chọn đề tài}

Trong quy trình đào tạo, giảng viên cần biên soạn câu hỏi, tạo nhiều đề,
quản lý lịch thi, theo dõi bài làm và tổng hợp kết quả. Khi các công việc
này thực hiện rời rạc, việc tái sử dụng học liệu và theo dõi trạng thái
kiểm tra trở nên khó khăn. AI ExamHub DAU được xây dựng như một nền tảng
thống nhất cho ngân hàng câu hỏi, hỗ trợ sinh câu hỏi từ tài liệu, quy trình
duyệt đề và tổ chức lớp học, phòng thi trực tuyến.

Đề tài phù hợp với học phần trí tuệ nhân tạo vì kết hợp một ứng dụng thực
tiễn với thành phần xử lý ngôn ngữ: đọc nội dung học liệu, đề xuất câu hỏi
và liên kết câu hỏi với bằng chứng trong nguồn. Hệ thống cũng có phương án
dự phòng dựa trên quy tắc khi không cấu hình dịch vụ mô hình ngôn ngữ. Việc
khảo sát đồng thời phần AI và quy trình web giúp làm rõ giới hạn của một
nguyên mẫu tích hợp, thay vì đồng nhất việc kết nối API với một mô hình đã
được huấn luyện và đánh giá độc lập.

\section{Mục tiêu của đồ án}

Đồ án hướng tới các mục tiêu sau:
\begin{itemize}[leftmargin=1.6em, itemsep=2pt]
  \item Xây dựng ứng dụng web gồm giao diện người dùng, API backend và cơ sở
        dữ liệu quan hệ.
  \item Hỗ trợ tải tài liệu học tập, trích xuất văn bản và đề xuất câu hỏi
        trắc nghiệm/tự luận có thông tin căn cứ.
  \item Hỗ trợ quản lý câu hỏi, tạo và duyệt đề thi, tổ chức phòng thi trực
        tuyến và lớp học theo vai trò.
  \item Kiểm tra các luồng chính bằng kiểm thử API và xác nhận giao diện có
        thể build/type-check.
  \item Nêu rõ các giới hạn của phần AI, dữ liệu và kiểm thử, không suy diễn
        kết quả chức năng thành độ chính xác mô hình.
\end{itemize}

\section{Câu hỏi nghiên cứu và tiêu chí thành công}

Phạm vi đồ án được định hướng bởi ba câu hỏi. Thứ nhất, làm sao kết nối việc
tạo đề, duyệt đề và tổ chức kiểm tra vào cùng một quy trình phần mềm có phân
quyền? Thứ hai, ở mức một nguyên mẫu, làm thế nào hỗ trợ khai thác tài liệu
học tập mà vẫn giữ được khả năng giảng viên rà soát nguồn? Thứ ba, những
bằng chứng nào là đủ để kết luận về độ hoàn thiện phần mềm, và những câu hỏi
nào cần một nghiên cứu đánh giá AI riêng?

Tiêu chí thành công của phần mềm là các chức năng trọng tâm có API/UI tương
ứng, quyền truy cập được phân theo vai trò, dữ liệu chính được lưu trong
PostgreSQL, các luồng nghiệp vụ có kiểm thử tự động và frontend build
thành công. Với thành phần AI, tiêu chí của nguyên mẫu chỉ là tạo được đầu
ra có cấu trúc hoặc nhánh dự phòng và cung cấp thông tin nguồn để kiểm tra.
Độ đúng, tính sư phạm, độ khó và hiệu quả học tập không được xem là đã xác
nhận cho đến khi có bộ chuẩn đánh giá độc lập.

\section{Phương pháp thực hiện}

Đồ án sử dụng phương pháp phát triển theo hướng tích hợp và kiểm chứng
nghiệp vụ. Trước hết khảo sát cấu trúc mã nguồn và schema, xác định các vai
trò cùng luồng sử dụng; tiếp đó mô tả kiến trúc client--server, mô hình dữ
liệu và hợp đồng API. Sau khi hiện thực frontend, backend và các migration,
các luồng được xác minh bằng kiểm thử API hiện có và lệnh build của
TypeScript/Vite. Cuối cùng, kết quả được phân loại thành bằng chứng chức
năng, bằng chứng build và các đánh giá chưa có dữ liệu hỗ trợ.

Phương pháp này phù hợp với mục tiêu tạo nguyên mẫu tích hợp trong phạm vi
đồ án, nhưng không thể thay thế thực nghiệm có người dùng, đo tải, kiểm thử
bảo mật độc lập hoặc đánh giá chất lượng nội dung AI bởi chuyên gia. Phạm
vi và loại bằng chứng được tóm tắt trong Bảng~\ref{tab:phuongphap}.

\begin{table}[H]
\centering
\caption{Phương pháp và bằng chứng tương ứng với câu hỏi của đồ án.}
\label{tab:phuongphap}
\begin{tabularx}{\textwidth}{>{\bfseries\color{daunavy}}p{0.25\textwidth} X X}
\toprule
\rowcolor{gray100}Câu hỏi & Cách khảo sát & Bằng chứng có/không có \\
\midrule
Tích hợp nghiệp vụ & Đọc frontend, route API, middleware và schema; chạy luồng kiểm thử. & Mã nguồn, 35 ca API đạt và build frontend thành công. \\
Hỗ trợ khai thác tài liệu & Đọc xử lý trích xuất, cấu hình mô hình, kiểm tra đầu ra và fallback. & Có kiểm thử chức năng; chưa có chấm điểm câu hỏi bởi chuyên gia. \\
Mức độ chất lượng AI & Xác định tập chuẩn và độ đo cần thiết. & Chưa có tập gán nhãn; không báo cáo accuracy/F1/độ đúng nội dung. \\
\bottomrule
\end{tabularx}
\end{table}

\section{Đóng góp và giới hạn đóng góp}

Đóng góp chính là tích hợp được quy trình tạo câu hỏi có căn cứ, chọn ngân
hàng để lập đề, duyệt và phát hành, tổ chức phòng thi, lưu tiến độ của
sinh viên và ghi nhận kết quả trên cùng nền tảng. Về kỹ thuật, hệ thống
minh hoạ cách phối hợp giao diện React/TypeScript, REST API Express và
PostgreSQL, có Docker Compose cho môi trường phát triển. Thành phần AI được
đóng góp ở mức tích hợp suy luận dịch vụ ngoài khi cấu hình cùng một nhánh
sinh dự phòng; đồ án không tuyên bố đề xuất kiến trúc mô hình mới, huấn
luyện mô hình riêng hay tạo ra một corpus học thuật mới.

\section{Phạm vi, đối tượng và giả định}

Đối tượng sử dụng gồm quản trị viên, giảng viên và sinh viên. Phạm vi triển
khai khảo sát trong báo cáo tập trung vào ứng dụng web, API Node.js/Express
và PostgreSQL. Luồng tạo câu hỏi nhận nội dung tài liệu PDF, DOC/DOCX hoặc
TXT; khả năng trích xuất phụ thuộc vào định dạng và chất lượng văn bản/OCR.
Mô hình ngôn ngữ bên ngoài chỉ được sử dụng khi nhà triển khai cung cấp cấu
hình cần thiết; nếu không, hệ thống dùng bộ sinh dự phòng dựa trên văn bản.

Đồ án không xây dựng hay huấn luyện mô hình nền tảng riêng, không có bộ dữ
liệu gán nhãn chuẩn để đo độ chính xác sinh câu hỏi, không triển khai các
thư mục AI riêng như những dịch vụ độc lập, và chưa đánh giá tải hoặc bảo
mật ở quy mô sản xuất. Chấm điểm tự động hiện tập trung vào câu trắc nghiệm;
bài tự luận cần giảng viên chấm theo nội dung và tiêu chí.

Các giả định gồm: người tạo đề có quyền sử dụng tài liệu được tải lên; nhà
triển khai tự cấu hình dịch vụ AI nếu muốn sử dụng nhánh mô hình ngôn ngữ;
giảng viên xem xét câu hỏi/đáp án trước khi phát hành; cơ sở dữ liệu được
khởi tạo theo schema/migration tương ứng. Những giả định này cần được rà
soát khi đưa hệ thống vào sử dụng tại một cơ sở đào tạo thật.

\section{Bố cục báo cáo}

Phần còn lại của báo cáo gồm ba chương. Chương~1 trình bày cơ sở lý thuyết
và công nghệ liên quan; Chương~2 phân tích yêu cầu, thiết kế dữ liệu và
kiến trúc, đồng thời mô tả phần cài đặt; Chương~3 tổng hợp kết quả, bằng
chứng kiểm thử, hạn chế và hướng phát triển. Sau phần tài liệu tham khảo là
phụ lục hướng dẫn chạy hệ thống, từ điển dữ liệu và danh mục API.
